\section*{Bab \ref{ch:b1:layered-earth} --- Bumi Berlapis}

\begin{solution}{exo:b1:layered-earth:1}
Fraksi volume adalah kubus dari perbandingan jari-jari: inti $(3480/6371)^3 = 0.1630$, yaitu \qty{16.3}{\percent}; \omterm{def:b1:layered-earth:layers}{inti dalam} $(1221.5/6371)^3 = 0.00705$, \qty{0.70}{\percent}.
\end{solution}

\begin{solution}{exo:b1:layered-earth:2}
Di atas: $\Phi = 13.7166^2 - \tfrac43\times7.26466^2 = 188.145 - 70.367 =
\qty{117.78}{km^2.s^{-2}}$, $K = \rho\Phi = \qty{5566.45}{kg.m^{-3}}\times
\qty{117.78e6}{m^2.s^{-2}} = \qty{655.6}{GPa}$. Di bawah: $\Phi = 8.06482^2 =
\qty{65.04}{km^2.s^{-2}}$, $K = 9903.49\times\num{65.04e6} = \qty{644.1}{GPa}$.
Modulus bulk hampir tidak berubah (\qty{-1.7}{\percent}). Yang berubah adalah modulus geser, $\mu = \rho v_S^2 = \qty{293.8}{GPa}$ di atas dan $0$ di bawah (cairan), dan massa jenis, yang naik sebesar \qty{78}{\percent}: karena itu $v_P$ turun meskipun $K$ hampir sama.
\end{solution}

\begin{solution}{exo:b1:layered-earth:3}
$\sin i_c = 6.2/8.1 = 0.7654$, $i_c = \ang{49.95}$. Jarak kritis
$2\times30\times\tan\ang{49.95} = 60\times1.1906 = \qty{71.4}{km}$. \omterm{def:b1:layered-earth:head-wave}{Jarak persilangan} $60\sqrt{14.3/1.9} = 60\times2.743 = \qty{164.6}{km}$.
\end{solution}

\begin{solution}{exo:b1:layered-earth:4}
$g = Gm/r^2 = \num{6.674e-11}\times\num{1.94e24}/(\num{3.48e6})^2 =
\qty{10.69}{m.s^{-2}}$, \qty{9}{\percent} lebih besar daripada di permukaan. Menuju inti, jari-jari kuadrat turun menjadi $(3480/6371)^2 = 0.298$ dari nilai permukaannya, sementara massa penarik hanya turun menjadi $1.94/5.972 = 0.325$ dari total, karena inti memuat sepertiga massa: $0.325/0.298 = 1.09$.
\end{solution}

\begin{solution}{exo:b1:layered-earth:5}
Dekat sumber $v_1 = 60/10.0 = \qty{6.0}{km.s^{-1}}$. Jauh dari sumber
$v_2 = (300 - 200)/(45.8 - 33.3) = 100/12.5 = \qty{8.0}{km.s^{-1}}$; waktu intersep adalah $t_i = 33.3 - 200/8.0 = \qty{8.3}{s}$. Lalu $\sin i_c =
0.75$, $\cos i_c = 0.6614$, dan $h = t_i v_1/(2\cos i_c) = 8.3\times6.0/1.3229
= \qty{37.6}{km}$. Periksa: $x_c = 2\times37.6\times\sqrt{14/2} = \qty{199}{km}$,
tempat kedua gelombang tiba pada $199/6.0 = 199/8.0 + 8.3 = \qty{33.2}{s}$.
\end{solution}

\begin{solution}{exo:b1:layered-earth:6}
Untuk $\rho$ seragam, $I = \tfrac{8\pi}{3}\rho\int_0^R r^4\,\mathrm dr =
\tfrac{8\pi}{15}\rho R^5$ dan $M = \tfrac43\pi\rho R^3$, sehingga
$I/MR^2 = (8/15)/(4/3) = 0.4$. Dua lapisan ($r_c = \qty{3.48e6}{m}$,
$R = \qty{6.371e6}{m}$): $M = \tfrac43\pi\bigl(11000\,r_c^3 + 4450\,(R^3 -
r_c^3)\bigr) = \qty{5.977e24}{kg}$, \qty{0.08}{\percent} di atas massa terukur;
$I = \tfrac{8\pi}{15}\bigl(11000\,r_c^5 + 4450\,(R^5 - r_c^5)\bigr)$ memberikan
$I/MR^2 = 0.3457$. Massanya sudah tepat, tetapi faktornya \qty{4.5}{\percent}
terlalu tinggi: \omterm{def:b1:layered-earth:layers}{mantel} nyata tidak seragam, massa jenisnya bertambah dengan kedalaman, yang memindahkan massa ke arah dalam.
\end{solution}

\begin{solution}{exo:b1:layered-earth:7}
$\Phi = 12.16^2 - \tfrac43\times6.673^2 = 147.87 - 59.37 =
\qty{88.49}{km^2.s^{-2}}$. Dengan kedalaman $z$,
$\mathrm d\rho/\mathrm dz = \rho g/\Phi = 4860\times9.9/\num{88.49e6} =
\qty{5.44e-4}{kg.m^{-4}}$, yaitu \qty{0.544}{kg.m^{-3}} per kilometer, atau
\qty{0.54}{g.cm^{-3}} per \qty{1000}{km}. PREM naik sebesar $5.357 - 4.844 =
0.513$ pada \qty{1000}{km} ini, \qty{6}{\percent} lebih kecil: gradien yang dihitung di puncak diterapkan pada seluruh interval, sementara $\Phi$ bertambah dengan kedalaman dan gradiennya menurun.
\end{solution}

\begin{solution}{exo:b1:layered-earth:8}
$P \approx \bar\rho g z = 3570\times9.9\times\num{670e3} = \qty{23.7}{GPa}$
(PREM: \qty{23.9}{GPa}). Itulah tekanan, sekitar \qty{24}{GPa}, tempat
ringwoodit terurai menjadi bridgmanit dan ferropericlase: diskontinuitas
\qty{660}{km}.
\end{solution}

\begin{solution}{exo:b1:layered-earth:9}
(a)~Besi cair: $10.90\times1.05 = 11.4$ hingga $10.90\times1.10 =
\qty{12.0}{g.cm^{-3}}$. (b)~$12.764/12.166 = 1.049$: lompatan
sebesar \qty{4.9}{\percent} dalam PREM, sedikit lebih besar daripada \qty{4.5}{\percent} dari
kajian terkini. (c)~Pembekuan saja hanya akan membuat padatan sedikit lebih padat
daripada cairan; lompatan yang lebih besar berarti padatan memuat lebih sedikit
unsur ringan daripada cairan: unsur-unsur itu lebih suka di cairan, dan pembekuan di
\omterm{def:b1:layered-earth:discontinuity}{batas inti dalam} meninggalkan sebagian besar unsur itu di belakang.
\end{solution}

\begin{solution}{exo:b1:layered-earth:10}
Misalkan $x_* = 2h\tan i_c$. Maka $\sqrt{x_*^2 + 4h^2} = 2h/\cos i_c$ dan
$t_r(x_*) = 2h/(v_1\cos i_c)$. \omterm{def:b1:layered-earth:head-wave}{Gelombang kepala} memberikan
$t(x_*) = 2h\tan i_c/v_2 + 2h\cos i_c/v_1 = \frac{2h}{v_1}\bigl(\sin^2 i_c/\cos
i_c + \cos i_c\bigr) = 2h/(v_1\cos i_c)$, dengan $\sin i_c = v_1/v_2$: sama.
Gradien hiperbola adalah $t_r'(x) = x/(v_1\sqrt{x^2 + 4h^2})$; pada $x_*$ ia adalah
$\sin i_c/v_1 = 1/v_2$, gradien garis \omterm{def:b1:layered-earth:head-wave}{gelombang kepala}. Karena
$t_r'' > 0$, hiperbola itu cembung dan berada di atas setiap garis singgungnya: untuk
$x > x_*$, $t_r > t_{\text{kepala}}$. Pada jarak kritis pantulan
menumbuk batas tepat pada $i_c$ dan merupakan kelahiran \omterm{def:b1:layered-earth:head-wave}{gelombang kepala}; di luar itu,
\omterm{def:b1:layered-earth:head-wave}{gelombang kepala}, yang diuntungkan karena merambat di medium cepat, tiba lebih dulu.
\end{solution}

\begin{solution}{exo:b1:layered-earth:11}
$m(r) = \tfrac43\pi\rho r^3$, sehingga $g = Gm/r^2 = \tfrac43\pi G\rho r$. Lalu
$P(r) = \int_r^R \rho\,\tfrac43\pi G\rho s\,\mathrm ds = \tfrac23\pi G\rho^2(R^2
- r^2)$. Di pusat $P = \tfrac23\pi\times\num{6.674e-11}\times5513^2\times
(\num{6.371e6})^2 = \qty{172}{GPa}$, \qty{47}{\percent} dari
\qty{363.85}{GPa} milik PREM: memusatkan massa ke arah pusat kira-kira menggandakan
tekanan pusat.
\end{solution}

\begin{solution}{exo:b1:layered-earth:12}
Dengan $z$ kedalaman, $\mathrm d\rho/\mathrm dz = \rho g/\Phi$; jika $g/\Phi$ adalah
konstanta $a$, $\mathrm d(\ln\rho)/\mathrm dz = a$ dan $\rho = \rho_0
e^{a(z - z_0)}$. Di sini $a(z - z_0) = 7.9\times\num{2258.5e3}/\num{87e6} =
0.2051$ dan $\rho = 9.90\times e^{0.2051} = 9.90\times1.2276 =
\qty{12.15}{g.cm^{-3}}$, \qty{0.2}{\percent} di bawah $12.17$ milik PREM: seluruh
kenaikan massa jenis melalui \omterm{def:b1:layered-earth:layers}{inti luar} adalah pemampatan diri.
\end{solution}

\begin{solution}{pb:b1:layered-earth:1}
\textbf{1.} $V = \tfrac43\pi(\num{6.371e6})^3 = \qty{1.0832e21}{m^3}$.

\textbf{2.} $\tfrac43\pi(r_2^3 - r_1^3)$: \omterm{def:b1:layered-earth:layers}{kerak}, \omterm{def:b1:layered-earth:layers}{mantel} atas dan zona transisi
\qty{3.0706e20}{m^3}; \omterm{def:b1:layered-earth:layers}{mantel} bawah \qty{5.9961e20}{m^3}; \omterm{def:b1:layered-earth:layers}{inti luar}
\qty{1.6890e20}{m^3}; \omterm{def:b1:layered-earth:layers}{inti dalam} \qty{7.634e18}{m^3}.

\textbf{3.} \qty{28.35}{\percent}, \qty{55.36}{\percent}, \qty{15.59}{\percent}
dan \qty{0.70}{\percent}: \omterm{def:b1:layered-earth:layers}{mantel} bawah adalah kulit terbesar, lebih dari separuh
volume Bumi.

\textbf{4.} $15.59 + 0.70 = \qty{16.3}{\percent}$.

\textbf{5.} \omterm{def:b1:layered-earth:layers}{Kerak}, \omterm{def:b1:layered-earth:layers}{mantel} atas dan zona transisi $\num{3.0706e20}\times3570 =
\qty{1.0962e24}{kg}$; \omterm{def:b1:layered-earth:layers}{mantel} bawah \qty{2.9381e24}{kg}; \omterm{def:b1:layered-earth:layers}{inti luar}
\qty{1.8410e24}{kg}; \omterm{def:b1:layered-earth:layers}{inti dalam} \qty{9.841e22}{kg}.

\textbf{6.} \qty{5.9737e24}{kg}, \qty{0.03}{\percent} di atas massa yang terukur
\qty{5.972e24}{kg}.

\textbf{7.} $\num{5.9737e24}/\num{1.0832e21} = \qty{5515}{kg.m^{-3}}$.

\textbf{8.} $\num{9.841e22}/\num{5.9737e24} = \qty{1.65}{\percent}$.

\textbf{9.} $(1.0962 + 2.9381)/5.9737 = \qty{67.5}{\percent}$.

\textbf{10.} Inti adalah \qty{16.3}{\percent} dari volume dan
\qty{32.5}{\percent} dari massa, dua kali lipat: massa jenis rata-ratanya,
\qty{10986}{kg.m^{-3}}, kira-kira $2.5$ kali massa jenis \omterm{def:b1:layered-earth:layers}{mantel}.

\textbf{11.} Kulit dengan massa jenis seragam memberikan
$I = \tfrac{8\pi}{3}\rho\int_{r_1}^{r_2} r^4\,\mathrm dr = \tfrac{8\pi}{15}\rho
(r_2^5 - r_1^5)$.

\textbf{12.} \omterm{def:b1:layered-earth:layers}{Kerak}, \omterm{def:b1:layered-earth:layers}{mantel} atas dan zona transisi \qty{2.676e37}{kg.m^2}; \omterm{def:b1:layered-earth:layers}{mantel} bawah
\qty{4.525e37}{kg.m^2}; \omterm{def:b1:layered-earth:layers}{inti luar} \qty{9.27e36}{kg.m^2}; \omterm{def:b1:layered-earth:layers}{inti dalam}
\qty{5.87e34}{kg.m^2}.

\textbf{13.} \omterm{def:b1:layered-earth:layers}{Mantel} bawah: ia memiliki separuh massa, pada jari-jari besar. \omterm{def:b1:layered-earth:layers}{Inti dalam}
memuat \qty{1.6}{\percent} dari massa, dan seluruhnya dekat sumbu: bagiannya dari momen adalah \qty{0.07}{\percent}.

\textbf{14.} $I = \qty{8.1345e37}{kg.m^2}$ dan $I/MR^2 = \num{8.1345e37}/
(\num{5.9737e24}\times(\num{6.371e6})^2) = 0.3355$.

\textbf{15.} \qty{1.4}{\percent} di atas $0.3308$, dan jauh di bawah $0.4$: model memiliki konsentrasi massa ke arah pusat yang hampir tepat.

\textbf{16.} Ya. Di dalam setiap kulit massa jenis nyata lebih tinggi di bagian bawah
daripada di atas; kulit seragam dengan massa yang sama menempatkan terlalu banyak massa itu
dekat jari-jari luarnya, tempat ia dihitung dengan $r^2$, dan karena itu melebihkan $I$.

\textbf{17.} $g = GM/R^2 = \num{6.674e-11}\times\num{5.9737e24}/
(\num{6.371e6})^2 = \qty{9.822}{m.s^{-2}}$.

\textbf{18.} Massa inti adalah $\num{1.8410e24} + \num{9.841e22} =
\qty{1.9394e24}{kg}$, dan $g = G m/r_c^2 = \qty{10.69}{m.s^{-2}}$.

\textbf{19.} $(\num{1.0962e24} + \num{2.9381e24})/(\num{3.0706e20} +
\num{5.9961e20}) = \num{4.0343e24}/\num{9.0667e20} = \qty{4450}{kg.m^{-3}}$.

\textbf{20.} $P \approx 4450\times10\times\num{2.891e6} = \qty{128.6}{GPa}$,
\qty{5}{\percent} di bawah \qty{135.75}{GPa} milik PREM.

\textbf{21.} $\rho_c = \num{1.9394e24}/\num{1.7653e20} = \qty{10986}{kg.m^{-3}}$,
dan $\Delta P = \tfrac23\pi\times\num{6.674e-11}\times10986^2\times
(\num{3.48e6})^2 = \qty{204.3}{GPa}$.

\textbf{22.} $128.6 + 204.3 = \qty{332.9}{GPa}$, \qty{8.5}{\percent} di bawah
\qty{363.85}{GPa} milik PREM. Kedua perkiraan itu mengabaikan kenaikan massa jenis terhadap kedalaman di dalam setiap wilayah, yang menempatkan lebih banyak massa lebih dalam, menjaga gravitasi tetap lebih tinggi dalam perjalanan ke bawah, dan menaikkan $\rho g$, gradien tekanan.

\textbf{23.} Massa total (soal 6), \omterm{def:b1:comparative-planetology:moi}{faktor momen inersia}
(soal 15) dan gravitasi permukaan (soal 17, terhadap rata-rata terukur
\qty{9.82}{m.s^{-2}}) diperiksa terhadap pengukuran Bumi; gravitasi di inti dan tekanan-tekanannya hanya dapat dibandingkan dengan PREM, yang sendiri adalah sebuah model.

\textbf{24.} Inti memuat \textbf{\qty[text-series-to-math=true]{32.5}{\percent}}
dari massa Bumi.
\end{solution}
